\documentclass[10pt,a4paper]{amsart}
\usepackage[margin=19mm]{geometry}
\usepackage{amsmath,amssymb,mathtools}
\usepackage[scr=rsfs,cal=euler]{mathalfa}
\usepackage{xfrac}
\usepackage[unicode,hidelinks,bookmarks=false]{hyperref}
\newcommand{\ie}{i.e.\ }
\newcommand{\cf}{cf.\ }
\newcommand{\resp}{respectively\ }
\newcommand{\Subsection}{\subsection}
\providecommand{\nobreakdash}{\nobreak\hbox{-}}
\setlength{\emergencystretch}{2em}
\multlinegap0pt
\usepackage{amssymb}
\usepackage{dsfont}
\newtheorem{lem}{Lemma}
\newcommand{\Ord}{{\rm Ord}}
\newcommand{\Z}{\mathbb Z}
\title{On the role of higher roots in prime ideal races}
\author{Mounir Hayani}
\date{Source: arXiv:2607.23150v1 (CC BY 4.0)}
\begin{document}
\maketitle

\noindent\emph{Source and licence.} On the role of higher roots in prime ideal races. Version arXiv:2607.23150v1 (CC BY 4.0), distributed by arXiv under the Creative Commons Attribution 4.0 International licence, http://creativecommons.org/licenses/by/4.0/, as recorded in the arXiv licence metadata for this exact version and shown on the abstract page https://arxiv.org/abs/2607.23150v1. Full licence notice: \texttt{licenses/arxiv-2607.23150-CC-BY-4.0.txt}.\par\smallskip

\noindent\textit{Editorial scope.} Counted unit: the article's lem difference2proots (source0.tex lines 733--735). The article's complete original proof of it is delivered with it (source0.tex lines 736--770, one \texttt{proof} environment of 35 lines). No mathematics is written, repaired or re-derived: the statement and the proof are the article's own lines, in the article's own notation, and no retained line is substituted or re-flowed.  Retained uncounted as the source-local context the proof needs: lines 514--517; lines 709--711; lines 713--715.  Nothing was omitted from any retained span, so the package carries no omission record.  No retained line contains a citation, so the wrapper carries no bibliography and no bibliography entry.  No retained line contains a figure environment, an image-inclusion command or drawing code, so no image is delivered and the package declares no asset dependency.  Editorial formatting only, in addition: an AMS \texttt{amsart} preamble that restates the article's own declarations and macros used by the retained lines, namely the environments \texttt{lem} on separate counters, together with the standard packages those lines need.  The retained environments print in this document as Lemma 1 in order of appearance, whereas the article prints the same items with the numbering of its own full document.  The complete native source of the article as distributed in the arXiv e-print for this exact version is bundled unchanged as \texttt{sources/206007/source0.tex}.
\begin{lem}\label{rootorder}
    Let $G$ be a finite group and let $n\ge 1$. If $g\in G$ is such that $d=\gcd(\Ord(g),n)$, then for all $x\in \mathcal{R}_n(g)$, we have \[ d\Ord(g) \mid \Ord(x)\mid n\Ord(g)\, .\]
    In other words, $\Ord(x)$ is a multiple of $d\Ord(g)$ that divides $n\Ord(g)$.
\end{lem}

    Let $G$ be a finite group and let $p\ge2$ be a prime number. For all $g\in G $ with order $\Ord(g)=p^rm$, such that $\gcd(p,m)=1$ and $r\ge0$, we have \begin{equation}\label{p-rootsnumber}
        r_p(g)=\#\{t\in \mathcal{R}_p\bigl(g^m\bigr)\, \colon\,  tg=gt\}\, .
    \end{equation}

    \begin{align}
        r_{2p}(g)=\#\left\{(s,t)\in \mathcal{R}_2\bigl(g\bigr) \times\mathcal{R}_p\bigl(g^m\bigr)\, \colon\,  st=ts\right\}\label{2p-rootsnumb2}\, ,
    \end{align}

\begin{lem}\label{difference2proots}
    Let $G$ be a finite group, let $p\ge 3$ be a prime number, and let $(x,y)\in G^2$ such that $\Ord(x)=\Ord(y)$. If $r_2(x)=r_2(y)$ and $r_p(x)=r_p(y)$, then $r_{2p}(x)-r_{2p}(y) \in p(p-1)\Z \cup p^2\Z$. If moreover $\Ord(x)$ and $\Ord(y)$ are even, then $r_{2p}(x)-r_{2p}(y) \in 2p(p-1)\Z\cup 2p^2\Z$.
\end{lem}

\begin{proof}
    For all $g\in G$ with order $\Ord(g)=p^\alpha \ell$ such that $\gcd(p,\ell)=1$, we define \[\mathcal{F}_p(g):=\left\{ t\in \mathcal{R}_p\bigl(g^\ell\bigr)\, \colon\,  tg=gt\right\}\, .\]
    From the equality $r_2(x)r_p(x)=r_2(y)r_p(y)$ and~\eqref{p-rootsnumber}, we deduce 
    \[\sum_{t\in \mathcal{F}_p(x)}\left| \mathcal{R}_2(x)\right|=\sum_{t\in\mathcal{F}_p(y)}\left|\mathcal{R}_2(y)\right|\, .\]
    This, combined with~\eqref{2p-rootsnumb2}, implies that
    \begin{align*}
        r_{2p}(x)-r_{2p}(y)&=\sum_{t\in \mathcal{F}_p(x)}\#\left\{s\in \mathcal{R}_2(x)\, \colon\,  st=ts\right\}-\sum_{t\in \mathcal{F}_p(y)}\#\left\{s\in \mathcal{R}_2(y)\, \colon\,  st=ts\right\}\\
        &=\sum_{t\in \mathcal{F}_p(y)}\#\left\{s\in \mathcal{R}_2(y)\, \colon\,  st\ne ts\right\}-\sum_{t\in \mathcal{F}_p(x)}\#\left\{s\in \mathcal{R}_2(x)\, \colon\,  st\ne ts\right\}
    \end{align*}
    Let $t\in \mathcal{F}_p(x)\setminus\{1\}$. We prove that we have $\#\left\{s\in \mathcal{R}_2(x)\, \colon\,  st\ne ts\right\}\in p\Z$ and if moreover $\Ord(x)$ is even, then $\#\left\{s\in \mathcal{R}_2(x)\, \colon\,  st\ne ts\right\}\in 2p\Z$.
    Assume that $\left\{s\in \mathcal{R}_2(x)\, \colon\,  st\ne ts\right\}\ne \emptyset$. Writing $\Ord(x)=p^rm$ with $\gcd(m,p)=1$, we have $t\in \mathcal{R}_{p}(x^m)$. Since $ts\ne st$ we have that $t\notin \langle s\rangle$, and thus $t\notin\langle x^m \rangle$. This means, by Lemma~\ref{rootorder}, that $\Ord(t)$ is a strict multiple of $p^{r}$ that divides $p^{r+1}$. Hence, $\Ord(t)=p^{r+1}$. We consider the action by conjugation of $\langle t\rangle$ on the set $\left\{s\in \mathcal{R}_2(x)\, \colon\,  st\ne ts\right\}$. This action has no trivial orbits, and the cardinality of each orbit divides $p^{r+1}$. Thus, $p$ divides $\#\left\{s\in \mathcal{R}_2(x)\, \colon\,  st\ne ts\right\}$. If $\Ord(x)$ is even, then by Lemma~\ref{rootorder}, for all $s\in \mathcal{R}_2(x)$, we have $\Ord(s)=2\Ord(x)$. Since $s^{\Ord(x)+1}\in \mathcal{R}_2(x)\setminus\{s\}$, we can write $\{s\in \mathcal{R}_2(x)\, \colon\,  st\ne ts\}$ as a disjoint union of subsets of the form $\left\{s,s^{\Ord(x)+1}\right\}$. This implies that $\#\{s\in \mathcal{R}_2(x)\, \colon\,  st\ne ts\}$ is a multiple of $2p$.
    
    Since $t\in \mathcal{R}_p(x^m)$, two cases arise: if $r=0$ (i.e., $\gcd(\Ord(x),p)=1$), then for all $h\in\langle t\rangle\setminus \{1\} $ we have $h\in \mathcal{F}_p(x)\setminus\{1\}$ and \[\left\{s\in \mathcal{R}_2(x)\, \colon\,  st\ne ts\right\}=\left\{s\in \mathcal{R}_2(x)\, \colon\,  sh\ne hs\right\}.\]
    Since there are $p-1$ elements in $\langle t\rangle\setminus\{1\}$, we have \[\sum_{h\in \langle t\rangle\setminus\{1\}} \#\{s\in \mathcal{R}_2(y)\, \colon\,  sh\ne hs\}\equiv\begin{cases}
        0\pmod{ p(p-1)}\quad&\text{if } \Ord(x)\text{ is odd}\\
        0\pmod{ 2p(p-1)} &\text{otherwise}.
    \end{cases}\]
    Since the latter holds for all $t\in \mathcal{F}_p(x)\setminus\{1\}$, we have that if $r=0$, then \[\sum_{t\in \mathcal{F}_p(x)}\#\left\{s\in \mathcal{R}_2(x)\, \colon\,  st\ne ts\right\}\equiv\begin{cases}
        0\pmod{ p(p-1)}\quad&\text{if } \Ord(x)\text{ is odd}\\
        0\pmod{ 2p(p-1)} &\text{otherwise}.
    \end{cases} \] 
    If $r\ge 1$, the set $\langle t\rangle\cap\mathcal{F}_p(x)=\langle t\rangle\cap\mathcal{R}_p(x^m)$ has cardinality $p$ (more generally, any element in a cyclic $p$-group has either no $p$-th roots or has exactly $p$ $p$-th roots). Similarly as above, for all these $p$ elements $h\in \langle t\rangle\cap\mathcal{F}_p(x) $ we have
    \[\left\{s\in \mathcal{R}_2(x)\, \colon\,  st\ne ts\right\}=\left\{s\in \mathcal{R}_2(x)\, \colon\,  sh\ne hs\right\}\, .\]
    Thus, if $r\ge1$, then
    \[\sum_{t\in \mathcal{F}_p(x)}\#\left\{s\in \mathcal{R}_2(x)\, \colon\,  st\ne ts\right\}=\begin{cases}
        0\pmod{ p^2}\quad&\text{if } \Ord(x)\text{ is odd}\\
        0\pmod{ 2p^2} &\text{otherwise}.
    \end{cases} \]  
    
    Similarly,  \[\sum_{t\in \mathcal{F}_p(y)}\#\left\{s\in \mathcal{R}_2(y)\, \colon\,  st\ne ts\right\}\in\begin{cases}
        p(p-1)\Z\cup p^2\Z\quad&\text{if } \Ord(y)\text{ is odd}\\
        2p(p-1)\Z\cup 2p^2\Z &\text{otherwise}.
    \end{cases} \]  
    This proves the lemma.
\end{proof}

\end{document}
